\section{Test Case 12: Platen stochastic GPU benchmark}

Test Case 12 combines the fixed 1,000-bead geometry and constant axial field
of Test Case 9 with the Platen integrator and stochastic air-drag parameters
of Test Case 4. Insertion, removal, evaporation, dynamic refinement, and
multiple-step Coulomb summation are disabled. The input is stored in
\texttt{examples/input-12/input.dat}.

Before timing begins, rank zero generates the complete Gaussian history in a
fixed step, bead, component, and draw order. For 1,000 steps the history
contains 6,006,000 double precision values, or 48,048,000 bytes. CPU
integration reads this layout directly and the OpenACC build copies it to the
accelerator once during initialization. Random values are neither generated
nor transferred during the temporal loop.
The stored history is limited to 100,000,000 double precision values (about
763 MiB). If a longer fixed run exceeds this capacity, indexing wraps to the
first stored timestep. This preserves CPU/GPU reproducibility but repeats the
noise periodically. Test Case 12 remains below the limit and does not wrap.

The persistent accelerator path performs the three Platen force evaluations,
the positive and negative predictors, stochastic velocity update, Heun
position and stress updates, and statistics on the device. An initial paired
measurement with NVFORTRAN 24.3 and one NVIDIA A30 required 29.050473 s on the
CPU and 1.874810 s with OpenACC, a speedup of 15.50. All six rows and fourteen
columns of the two sampled outputs were identical at the written precision.
The versioned results and comparison helper are stored in
\texttt{tests/performance/integrators/}.
